\documentclass[11pt,reqno]{amsart}
\usepackage{amsmath,amssymb,graphicx,amscd}
\usepackage{mathrsfs}
\usepackage{enumerate}
% pour les S des spheres
\usepackage{mathrsfs}
\usepackage{dsfont}
\usepackage{stmaryrd}
\usepackage{wrapfig}
\usepackage{pstricks}

% ----------------------------------------------------------------
\vfuzz2pt % Don't report over-full v-boxes if over-edge is small
\hfuzz2pt % Don't report over-full h-boxes if over-edge is small

\newcommand{\prob}{{\mathbb{P}}}

\renewcommand{\S}{\mathscr{S}}  % complex unit sphere

\newcommand{\CC}{{\mathbb C}}
\newcommand{\HH}{{\mathbb H}}
\newcommand{\RR}{{\mathbb R}}
\newcommand{\QQ}{{\mathbb{Q}}}
\newcommand{\NN}{{\mathbb{N}}}
\newcommand{\ZZ}{{\mathbb{Z}}}
\renewcommand{\P}{{\mathbb{P}}}
\renewcommand{\span}{{\mathrm{span}}}
\newcommand{\sgn}{\operatorname{{sgn}}}
 \DeclareMathOperator{\Prob}{\mathbb{P}}   %probability
 \DeclareMathOperator{\E}{\mathbb{E}}      %expectation
 \DeclareMathOperator{\cov}{cov}             %logarithmic integral
\newcommand{\1}{\mathds{1}}

 \newcommand{\I}{1\!\!1}                   %Indicator function


 \newcommand{\rd}{{\mathrm{d}}}            %differential
 \newcommand{\ii}{{\mathrm{i}}}
  \newcommand{\Id}{{\mathrm{Id}}}
\newcommand{\law}{\overset{\mbox{\rm \scriptsize law}}{=}}
\newcommand{\convlaw}{\overset{\mbox{\rm \scriptsize law}}{\longrightarrow}}


\renewcommand{\Re}{{\mathfrak{Re}}}
\renewcommand{\Im}{{\mathfrak{Im}}}

% THEOREMS -------------------------------------------------------

\newtheorem{thm}{Theorem}[section]
\newtheorem{cor}[thm]{Corollary}
\newtheorem{lem}[thm]{Lemma}
\newtheorem{prop}[thm]{Proposition}
\theoremstyle{definition}
\newtheorem{defn}[thm]{Definition}
\theoremstyle{remark}
\newtheorem*{rem}{Remark}



\numberwithin{equation}{section}


\begin{document}

~\vspace{-3cm}



\begin{center}
{\bf Probability, homework 7, due October 28.}
\end{center}
\vspace{0.5cm}
\small


\noindent{\bf Exercise 1}.
Read carefully Section 3.8 in Varadhan's {\it Probability Theory}.  Explain the proof of Theotrem 3.21 or Lemma 3.22 to a friend, on the board.\\






{\noindent \bf Long problem.} The goal of this problem is to prove the iterated logarithm law, first for Gaussian random variables. In other words, for $X_1, X_2\dots$ i.i.d. standard Gaussian random variables, denoting $S_n=X_1+\dots+X_n$, we have
\begin{equation}\label{eqn:ILL}
\mathbb{P}\left(
\limsup_{n\to\infty}\frac{S_n}{\sqrt{2n\log\log n}}=1
\right)
=1
\end{equation}


\begin{enumerate}[(1)]
\item Prove that
$$
\mathbb{P}(X_1>\lambda)\underset{\lambda\to\infty}{\sim}\frac{1}{\lambda\sqrt{2\pi}}e^{-\frac{\lambda^2}{2}}.
$$
In the following questions we denote $f(n)=\sqrt{2n\log\log n}$, $\lambda>1$, $c,\alpha>0$, $A_k=\{S_{\lfloor \lambda^k\rfloor}\geq c f(\lambda^k)\}$, 
$C_k=\{S_{\lfloor \lambda^{k+1}\rfloor}-S_{\lfloor \lambda^{k}\rfloor}\geq c f(\lambda^{k+1}-\lambda^k)\}$ and 
$D_k=\{\sup_{n\in\llbracket\lambda^k,\lambda^{k+1}\rrbracket}\frac{S_{n}-S_{\lfloor \lambda^{k}\rfloor}}{f(\lambda^k)}\geq \alpha\}$. 
\item Prove that for any $c>1$ we have $\sum_{k\geq 1}\mathbb{P}(A_k)<\infty$ and 
$$
\limsup_{k\to\infty}\frac{S_{\lfloor \lambda^k\rfloor}}{f(\lambda^k)}\leq 1\ \mbox{a.s.}
$$ 
\item Prove that for any $c<1$ we have $\sum_{k\geq 1}\mathbb{P}(C_k)=\infty$ and 
$$
\mathbb{P}(C_k\ {\rm i.o.})=1.
$$ 
\item Let $\varepsilon>0$ and choose $c=1-\varepsilon/10$. Prove that almost surely the following inequality holds for infinitely many $k$:
$$
\frac{S_{\lfloor \lambda^{k+1}\rfloor}}{f(\lambda^{k+1})}\geq c \frac{f(\lambda^{k+1}-\lambda^k)}{f(\lambda^{k+1})}-(1+\varepsilon)\frac{f(\lambda^k)}{f(\lambda^{k+1})}.
$$ 
\item By choosing a large enough $\lambda$ in the previous inequality, prove that almost surely 
$$
\limsup_{n\to\infty}\frac{S_n}{f(n)}\geq1.
$$
\item Prove that for any $n\in\llbracket\lambda^k,\lambda^{k+1}\rrbracket$ and $S_n>0$ we have
$$
\frac{S_n}{f(n)}\leq \frac{S_{\lfloor \lambda^k\rfloor}}{f(\lfloor\lambda^k\rfloor)}+
\frac{S_n-S_{\lfloor \lambda^k\rfloor}}{f(\lfloor\lambda^k\rfloor)}.
$$
\item Prove that
$$
\mathbb{P}(D_k)\underset{k\to\infty}{\sim}2\mathbb{P}\left(X_1\geq \frac{\alpha f(\lambda^k)}{\sqrt{\lambda^{k+1}-\lambda^k}}\right)
\underset{k\to\infty}{\sim}
\frac{{\rm cst}}{\sqrt{\log k}}\left(\frac{1}{k}\right)^{\frac{\alpha^2}{\lambda-1}}.
$$
\item Prove that for $\alpha^2>\lambda-1$, almost surely
$$
\limsup_{n\to\infty}\frac{S_n}{f(n)}\leq \limsup_{n\to\infty}\frac{S_{\lfloor\lambda^k\rfloor}}{f(\lambda^k)}+\alpha.
$$
\item  By choosing appropriate $\lambda$ and $\alpha$, prove that
 almost surely 
$$
\limsup_{n\to\infty}\frac{S_n}{f(n)}\leq1.
$$
\item State a result similar to (\ref{eqn:ILL}) for i.i.d. uniformly bounded random variables. Which steps in the above proof need to be modified to prove this universality result? How?
\end{enumerate}


\end{document} 